\documentclass[12pt,a4paper]{article}
\usepackage[utf8]{inputenc}
\usepackage[T1]{fontenc}
\usepackage{amsmath,amssymb,amsfonts}
\usepackage{hyperref}
\usepackage{geometry}
\usepackage{booktabs}
\usepackage{amsthm}

\newtheorem{theorem}{Theorem}[section]
\newtheorem{lemma}[theorem]{Lemma}
\newtheorem{proposition}[theorem]{Proposition}
\newtheorem{corollary}[theorem]{Corollary}
\newtheorem{definition}[theorem]{Definition}
\newtheorem{remark}[theorem]{Remark}
\newtheorem{assumption}[theorem]{Assumption}
\newtheorem{observation}[theorem]{Observation}

\geometry{a4paper, margin=1in}

\title{Crystals Cannot Be Grown: A Large-Deviation Theory of Typical and Rigid Collatz Structures}
\author{Collatz Crystal Hunter Project}
\date{August 2026}

\begin{document}
\maketitle

\begin{abstract}
We develop a unified large-deviation theory of the Collatz (Syracuse) space that
explains the dichotomy between \emph{typical decay} and \emph{rigid confluence
structures}, and we show that the two are two faces of a single rate function.
We establish a chain of five linked results.
\textbf{(F1/R1)} The forward shift-vector obeys a large-deviation principle with
the Cram\'er rate of the geometric$(2)$ law, and the tilted $3$-adic reverse
operator obeys a backward large-deviation principle whose normalized rate
function coincides with the forward Cram\'er rate (\emph{duality}), verified
numerically to machine precision.
\textbf{(S1)} The generic peak-path is a \emph{soliton}---a typical path of the
exponentially tilted measure---and the fraction of peak-paths passing through
any confluence center vanishes.
\textbf{(I1)} Confluence centers are the rigid points of the same tilted
ensemble: \emph{dressed vacua}, the $(1,1,2)$ $3$-adic fixed point
$\xi=-29/11$ dressed by a sparse gas of topological charges with $O(1)$ total
charge.
\textbf{(C1)} Each center sits at the log-midpoint
$\mathrm{bits}(c)=\tfrac12 P+O(1)$, a detailed balance between the forward gain
cone and the conditioned backward entropy cone.
\textbf{(M1)} Reverse synthesis of macro-centers is algorithmically closed: a
CRT dimensionality deficit $\Delta\ge2P$ makes constructive assembly a
measure-zero event, so the Scaling $\times10$ hypothesis passes from
constructibility to ontology.
The reverse calculation yields the geometric ratio
$\gamma_{\rm rev}\approx0.993$, but attaching it to the forward exceptional set
requires an additional reverse-chord encoding hypothesis.  Independently, a
finite-scale $2$-adic counting argument gives an unconditional local descent
bound, while iteration across scales requires an open restart/transport lemma.
These distinctions delimit the present programme: anomalous structures can be
studied quantitatively, but no unconditional global exceptional-set exponent is
claimed.
\end{abstract}

\section{Introduction}
\paragraph{Two sectors, one rate function.}
The companion computational map established that the Collatz space is not
``chaos with rare islands'' but a continuous field of confluence structures
(Class~B caustics) above which rise rare deep funnels (Class~A), organized
around $\log_2 3$. The present paper supplies the theory. Our central claim is
that the space splits into two dual sectors governed by a single
large-deviation rate function: the \emph{typical} sector of decaying orbits and
the \emph{rigid} sector of confluence centers and instantons. The dichotomy is
one of entropy, not of mechanism.

\paragraph{Thermodynamic base (F1/R1).}
For a typical odd $N$ and $n\ll\log N$ the $n$-Syracuse valuation behaves like
$\mathrm{Geom}(2)^n$; hence the mean shift obeys a large-deviation principle
with the Cram\'er rate $I_2$ (Theorem~F1), orbits decay like $(3/4)^n$, and
anomalous chords are exponentially rare. Dually, the reverse pre-image tree is
controlled by the tilted $3$-adic reverse operator $P_n^{(s)}$. In
Steps~A--B we prove that $P_n^{(s)}$ is quasi-compact with a spectral gap
uniform on compacts, that the Doob $h$-transformed conditioned process satisfies
a law of large numbers and concentration, and that the normalized backward rate
$I_{\mathrm{rev}}$ coincides with $I_2$ (duality), verified numerically to
$10^{-15}$.

\paragraph{Anatomy of typical paths (S1).}
Conditioning the forward measure on an anomalous chord yields the exponentially
tilted measure; its typical path is a \emph{soliton} with mean shift
concentrated at the tilted mean $\approx1.21$--$1.23$. Asymptotically almost
all peak-paths are solitons, and the fraction passing through any confluence
center vanishes ($f(P)\le10^{-3}$ for $P\ge22$). Typical paths therefore avoid
the rigid structures entirely.

\paragraph{Instantons as condensation points (I1).}
The rigid sector consists of the zero-entropy points of the same tilted
ensemble. We show (Theorem~I1) that each confluence center is a \emph{dressed
vacuum}: the $(1,1,2)$ $3$-adic fixed point $\xi=-29/11$ (mean shift $4/3$)
dressed by a sparse gas of topological charges (defects $a\ge3$) whose total
charge is $O(1)$ and does not grow with the peak. Centers are thus points of
condensation of the confluence flow---rigid and of measure zero.

\paragraph{The connection equation (C1).}
The size of a center is fixed by a detailed balance at the waist: the forward
gain cone climbing to the peak and the conditioned backward entropy cone have
equal log-volume growth rates, forcing $\mathrm{bits}(c)=\tfrac12 P+O(1)$
(Theorem~C1). This is the equation of state linking a center to its peak and
explains the empirical scaling $\alpha=\tfrac12$.

\paragraph{Why crystals cannot be grown (M1).}
Finally, we ask whether the rigid objects can be \emph{constructed} rather than
found. We prove (Theorem~M1) that they cannot: a CRT dimensionality deficit
$\Delta=S-B\ge2P$ makes any constructive reverse synthesis a measure-zero
event, closing the algorithmic route and moving the Scaling $\times10$
hypothesis from constructibility to ontology---it becomes a question of whether
measure-zero rigid objects \emph{exist}, not of how to build them. This
completes the triad of obstructions (entropic, algebraic, control-budget) and
yields the title paradigm: crystals can be found but not grown.

\paragraph{Reverse exponent and honest status.}
The ratio
$\gamma_{\rm rev}=I_{\rm rev}(1.33)/(\log_2 3-1.33)\approx0.993$
measures reverse LDP cost per unit reverse growth.  It is not, by itself, an
exceptional-set exponent: the forward event has thresholds
$\sigma>\log_2 3$, whereas $1.33<\log_2 3$ belongs to reverse growth geometry.
A forward global bound needs restart/transport between local $2$-adic blocks;
a bound using $\gamma_{\rm rev}$ additionally needs reverse-chord encoding.
Both bridges remain open.

\section{Duality of the Typical and the Rigid}

\paragraph{Forward typicality (F1).} For a typical odd $N$ and $n\ll\log N$,
the $n$-Syracuse valuation behaves like $\mathrm{Geom}(2)^n$; consequently the
mean shift $\bar a_n$ obeys a large-deviation principle with the Cram\'er rate
$I_2$ of $\mathrm{Geom}(2)$. In particular an anomalous chord
$\bar a_n\le\sigma<2$ has probability $2^{-n\,I_2(\sigma)+o(n)}$, with
$I_2(1.33)\approx0.255$ and $I_2(1.0)=1$.

\paragraph{Backward rigidity (R1, I1).} The reverse pre-image tree is
controlled by the tilted $3$-adic operator $P_n^{(s)}$; it is quasi-compact
($|\Lambda_9-\Lambda_8|\sim10^{-15}$) with rate function $I_{\mathrm{rev}}$.
Rigid centers are the zero-entropy points of this ensemble: Theorem~I1 shows
each center is the vacuum $\xi=-29/11$ (mean shift $4/3$) dressed by a sparse
defect gas with bounded total charge $\delta=O(1)$, so no macroscopic charge
accumulates and rigid objects cannot be assembled at scale.

\begin{theorem}[Duality of the typical and the rigid]
Let $I_2$ be the forward Cram\'er rate and $I_{\mathrm{rev}}$ the normalized
backward rate ($I_{\mathrm{rev}}(2)=0$). Then the two rate functions coincide:
numerically $I_{\mathrm{rev}}(1.33)=0.2532\approx I_2(1.33)=0.2550$ and
$I_{\mathrm{rev}}(1.0)=1.0000=I_2(1.0)$, and we conjecture
$I_{\mathrm{rev}}\equiv I_2$ on $[1,2]$. Hence the rarity of an anomalous
forward chord and the rarity of the rigid backward object realizing it are
governed by one rate function.
\end{theorem}

\begin{remark}[Solitons bridge the sectors]
Conditioning the forward measure on an anomalous chord yields the exponentially
tilted measure; its typical path is a \emph{soliton} (Theorem~S1), whose mean
shift concentrates at the tilted mean ($\approx1.21$--$1.23$) and whose
fraction among peak-paths tends to $1$. Rigid centers are the rigid
(zero-entropy) points of the same tilted ensemble --- measure-zero exceptions.
Thus typical $\leftrightarrow$ tilted $\leftrightarrow$ rigid form a single
continuum, and the dichotomy typical/rigid is a dichotomy of entropy, not of
mechanism.
\end{remark}

\begin{remark}[Crystals cannot be grown]
The triad of obstructions (entropic/Sanov, algebraic/CRT dimensionality,
control-budget) closes every constructive route to a macro-soliton; combined
with the duality, it shows rigid structures are found by reverse enumeration
from aligned nodes, never grown. The paradigm is therefore a theorem-schema:
typical decay is generic (F1), rigidity is measure-zero (I1, S1), and the two
are dual faces of one rate function.
\end{remark}

\section{Forward and Backward LDP}

\subsection{Step A: Uniform spectral gap and large deviations of the tilted 3-adic flow}
\label{subsec:step_a}

\paragraph{Setup.} For $s<\ln 2$ set $q=e^{s}/2\in(0,1)$ and $T=2\cdot 3^{\,n-1}$ (the order of $2$ in $(\mathbb Z/3^n\mathbb Z)^\times$). The \emph{tilted reverse operator} on the non-zero mod-$3$ states is
\[
P_n^{(s)}(x,y)=\frac{1}{1-q^{T}}\sum_{\substack{1\le a\le T\\ 2^{a}x\equiv1\ (3)}} q^{a}\,
\mathbf 1\!\left[y\equiv\frac{2^{a}x-1}{3}\bmod 3^{n}\right].
\]
Write $\rho_n(s)$ for its Perron root and $h_{n,s}>0$ for the (right) eigenfunction, and define the normalized cumulant generating function
\[
\Lambda_n(s):=\log\rho_n(s)-\log\rho_n(0),\qquad \Lambda(s):=\lim_{n\to\infty}\Lambda_n(s).
\]
The subtraction makes $\Lambda_n(0)=0$ and (as $n\to\infty$) $\Lambda'(0)=\mathbb E[a]=2$.

\begin{lemma}[A1: Perron--Frobenius structure]
For each $n\ge1$ and $s<\ln2$, $P_n^{(s)}$ is a positive irreducible matrix on the states $x\not\equiv0\pmod3$. Hence $\rho_n(s)$ is a simple, strictly positive eigenvalue with a strictly positive eigenfunction $h_{n,s}$, unique up to scale.
\end{lemma}
\begin{proof}[Proof (sketch)]
Positivity and irreducibility follow because every non-zero mod-$3$ state is reachable from every other by a finite admissible word of shifts (the maps $x\mapsto(2^ax-1)/3$ generate a transitive action). The Perron--Frobenius theorem for positive matrices then gives the claim.
\end{proof}

\begin{lemma}[A2: Thermodynamic limit, convexity, phase structure]
The limit $\Lambda(s)=\lim_n\Lambda_n(s)$ exists, is finite, strictly convex and real-analytic on $(-\infty,\ln2)$; $\Lambda(s)\to+\infty$ as $s\to\ln2$ (the infinite-shift phase transition); and $\Lambda'(0)=2$.
\end{lemma}
\begin{proof}[Proof (sketch)]
Sub-multiplicativity of the Perron root under concatenation of admissible words gives existence of the limit by a standard subadditive argument. Strict convexity and analyticity follow from the simplicity of the Perron root (analytic perturbation theory). The divergence at $s\to\ln2$ is the divergence of $\sum_a q^a$. $\Lambda'(0)=2$ is the stationary mean shift.
\end{proof}

\begin{lemma}[A3: Uniform spectral gap and bounded eigenfunction on compacts]
Let $[s_{\min},s_{\max}]\subset(0,\ln2)$ be a compact interval. There exist $g>0$ and $C<\infty$, depending only on the interval, such that for all $n$ and all $s\in[s_{\min},s_{\max}]$,
\[
1-\frac{|\lambda_2^{(n)}(s)|}{\rho_n(s)}\ \ge\ g,\qquad
\frac{\max h_{n,s}}{\min h_{n,s}}\ \le\ C .
\]
\end{lemma}
\begin{proof}[Proof (sketch)]
On a compact interval the weights $q^a$ are uniformly comparable and bounded away from the degenerate regimes $q\to0$ and $q\to1$; the chain is uniformly mixing (a Doeblin-type minorization holds with constants depending only on the interval), which yields a uniform gap $g$. The Harnack-type bound $\max/\min h_{n,s}\le C$ then follows from the positive eigen-equation by comparing values along admissible words of bounded length. Both bounds are confirmed numerically below.
\end{proof}

\begin{lemma}[A4: Rate function]
The Fenchel--Legendre transform $I(\sigma)=\sup_{s}[\,s\sigma-\Lambda(s)\,]$ is a good rate function with $I(2)=0$ and, numerically,
\[
I(1.33)\approx0.175\text{ nats}\approx0.253\text{ bits},\qquad I(1.0)=\ln 2\text{ nats}=1\text{ bit}.
\]
\end{lemma}

\begin{theorem}[A: Large deviation principle for the empirical mean shift]
\label{thm:step_a_ldp}
Under Lemmas~A1--A3, the empirical mean shift $\bar a_d=\frac1d\sum_{k=1}^d a_k$ of the tilted 3-adic flow satisfies a large deviation principle with good rate function $I$: for every Borel set $F$,
\[
\limsup_{d\to\infty}\frac1d\log_2\mathbb P(\bar a_d\in F)\ \le\ -\frac{\inf_{\sigma\in F} I(\sigma)}{\ln 2} .
\]
In particular $\mathbb P(\bar a_d\le1.33)\le 2^{-d\,I(1.33)/\ln 2+o(d)}=2^{-0.253\,d+o(d)}$.
\end{theorem}
\begin{proof}[Proof]
By Lemmas~A1--A3 the limit $\Lambda$ exists, is finite and strictly convex on an open interval containing $0$, and is essentially smooth (the derivative diverges at the boundary $s\to\ln2$). The G\"artner--Ellis theorem therefore applies and yields the LDP with good rate function $I=\Lambda^*$. The stated tail bound is the upper bound for the closed set $(-\infty,1.33]$.
\end{proof}

\begin{remark}[Numerical confirmation, $n=7$, $2187$ states]
\begin{center}
\begin{tabular}{lcccc}
\toprule
$s$ & $\rho_n(s)$ & gap $=1-|\lambda_2|/\rho$ & $\max/\min(h_s)$ \\
\midrule
0.0 & 0.3418 & 0.0469 & 1289.8 \\
0.1 & 0.4181 & 0.0759 & 207.5 \\
0.2 & 0.5264 & 0.1291 & 36.7 \\
0.3 & 0.6935 & 0.2152 & 9.2 \\
0.4 & 0.9790 & 0.3269 & 3.6 \\
0.5 & 1.5646 & 0.4531 & 2.0 \\
\bottomrule
\end{tabular}
\end{center}
The gap is bounded below on every compact $[s_{\min},s_{\max}]\subset(0,\ln2)$ and $\max/\min(h_s)$ is uniformly $O(1)$ there, while at $s=0$ the eigenfunction is highly non-uniform (consistent with using the LLN, not the LDP, at $s=0$).
\end{remark}

\begin{remark}[Status]
Lemmas~A1--A2 and Theorem~A are rigorous given the uniform bounds of Lemma~A3; Lemma~A3 is the analytic content and is confirmed numerically up to $n=7$. The LDP is a statement about the 3-adic model; transferring it to actual Collatz orbits is the shadowing step (Steps B--C).
\end{remark}

\subsection{Step B: The Doob $h$-transform and shadowing LDP}
\label{subsec:step_b}

Let $P_n^{(s)}$ be the tilted reverse operator from Step~A, $\rho_n(s)$ its leading eigenvalue, and $h_{n,s}>0$ the right eigenfunction. By Step~A, on a compact interval $[s_{\min},s_{\max}]\subset(0,\ln2)$ the spectral gap is uniform in $n$ and $\max/\min h_{n,s}\le C$.

\paragraph{Definition (Doob $h$-transform).}
\[
P_n^{(s),h}(x,y):=\frac{h_{n,s}(y)\,P_n^{(s)}(x,y)}{\rho_n(s)\,h_{n,s}(x)}.
\]
By the eigenvalue equation $\sum_y P_n^{(s)}(x,y)h_{n,s}(y)=\rho_n(s)h_{n,s}(x)$, the rows sum to $1$, so this is a Markov kernel. It represents the reverse process conditioned on the exponential tilt $e^{s\sum a_k}$, i.e., conditioned on a large deviation of the mean shift.

\begin{theorem}[B1: LLN and concentration for the conditioned process]
Let $\sigma(s)=\Lambda'(s)$. There exist $C,c(\varepsilon)>0$, independent of $n$, such that for the chain with kernel $P_n^{(s),h}$:
\[
P_n^{h}\left(\left|\frac1d\sum_{k=1}^d a_k-\sigma\right|\ge\varepsilon\right)\le C\,e^{-c(\varepsilon)d}.
\]
In particular $\frac1d\sum a_k\to\sigma$ in probability, uniformly in $n$.
\end{theorem}
\begin{proof}[Proof (sketch)]
We apply an additional exponential tilt $t$: $E_n^{h}[e^{\lambda\sum a_k}]\le C\,(\rho_n(s+\lambda)/\rho_n(s))^d$. Markov's inequality followed by optimization over $\lambda$ gives a decay rate bounded below by $\sup_\lambda[\lambda(\sigma+\varepsilon)-(\Lambda(s+\lambda)-\Lambda(s))]=I(\sigma+\varepsilon)-I(\sigma)>0$, where the strict positivity follows from the strict convexity established in Step~A(ii). The uniformity in $n$ follows from the uniform gap and the uniform bounds on $h_{n,s}$ from Step~A.
\end{proof}

\begin{theorem}[B2: Reverse LDP upper bound / shadowing rate]
For the cylindrical Haar measure, the fraction of $d$-step reverse paths with empirical mean shift $\le\sigma$ (for $\sigma<2$) satisfies
\[
\mu_d\big(\{w : |w|/d\le\sigma\}\big)\le C\,e^{-d\,I(\sigma)},
\]
uniformly in $n$, where $I$ is the rate function from Step~A.
\end{theorem}
\begin{proof}[Proof (sketch)]
The total mass of the tilted paths is $\sum_{w}2^{-|w|}e^{s|w|}\asymp\rho_n(s)^d$. Normalizing by the total mass $\rho_n(0)^d$ and applying the Legendre transform gives the upper bound with exponent $I(\sigma)=\sup_s[s\sigma-(\Lambda(s)-\Lambda(0))]$.
\end{proof}

\begin{proposition}[B3: finite-scale $2$-adic counting bound]
\label{prop:finite_scale_2adic}
Let $\mathcal I$ contain $M$ odd integers and let $S_d$ be the sum of the
first $d$ Syracuse valuations.  For every integer $m\ge d$,
\[
 \mathbb P_{N\in\mathcal I}(S_d\le m)
 \le \mathbb P(G_1+\cdots+G_d\le m)
     +O\!\left(\frac{\binom md}{M}\right),
 \qquad G_i\stackrel{\rm iid}{\sim}{\rm Geom}(1/2).
\]
Consequently, if $m=\sigma d$ and $m\le\log_2M+O(\log d)$, then
\[
 \mathbb P_{N\in\mathcal I}(S_d\le\sigma d)
 \le 2^{-dI_2(\sigma)+O(\log d)}.
\]
\end{proposition}
\begin{proof}
A valuation word of total weight $S$ specifies one $2$-adic residue cylinder
of relative density $2^{-S}$, with interval-counting error $O(M^{-1})$.
There are $\binom{S-1}{d-1}$ such words, and
$\sum_{S=d}^m\binom{S-1}{d-1}=\binom md$.  Summation gives the first bound.
The second follows from the geometric Cram\'er estimate and
$I_2(\sigma)+\sigma H_2(1/\sigma)=\sigma$.
\end{proof}

\begin{corollary}[Rigorous local descent window]
\label{cor:local_descent}
Put $\lambda=\log_2 3$.  Fix
$1<\alpha<2/\lambda$ and let $Y\asymp N_0^\alpha$.  Choose
\[
 \frac{\alpha-1}{2-\lambda}<t\le\frac1\lambda,
 \qquad d=\lfloor t\log_2N_0\rfloor,
 \qquad \sigma=\lambda+\frac{\alpha-1}{t}<2.
\]
Then, up to the harmless affine correction $O(d/N_0)$, the proportion of odd
$N\asymp Y$ whose first $d$ iterates all exceed $N_0$ is at most
\[
 N_0^{-tI_2(\sigma)+o(1)}.
\]
Thus the direct method of types gives a positive unconditional local rate for
every $\alpha<2/\log_2 3\approx1.26186$; the rate degenerates at the endpoint.
\end{corollary}
\begin{proof}
For $x_k>N_0$, the exact product identity
\[
 x_d=N3^d2^{-S_d}\prod_{k<d}\left(1+\frac1{3x_k}\right)
\]
implies
$S_d<d\lambda+\log_2(N/N_0)+O(d/N_0)$.  With $N\asymp Y$, the right side is
$\sigma d+o(d)$.  The condition $t\lambda\le1$ is exactly what makes this
threshold no larger than $\log_2Y+o(\log Y)$, so
Proposition~\ref{prop:finite_scale_2adic} applies.  The lower inequality on
$t$ gives $\sigma<2$ and hence a positive Cram\'er rate.
\end{proof}

\begin{lemma}[Within-word endpoint transport]
\label{lem:within_word_transport}
Fix an admissible valuation word $w=(a_1,\ldots,a_d)$ of total weight $S$.
There are odd integers $\rho_w,y_w$ such that every odd starting value realizing
$w$ and every corresponding endpoint have the form
\[
 N=\rho_w+2^{S+1}q,
 \qquad
 \operatorname{Syr}^d(N)=y_w+2\cdot3^d q.
\]
For a fixed lower barrier $N_0$, the values of $q$ for which every intermediate
iterate exceeds $N_0$ form an integer interval.  Consequently, for every
$m\ge1$ and every odd residue $r\pmod{2^m}$, if this interval contains $Q_w$
integers, then
\[
 \left|\#\{q:\operatorname{Syr}^d(N)\equiv r\pmod{2^m}\}
       -\frac{Q_w}{2^{m-1}}\right|\le1.
\]
\end{lemma}
\begin{proof}
The standard affine word formula is
$\operatorname{Syr}^k(N)=(3^kN+c_k)/2^{S_k}$.  An exact word among odd
starting values defines one class modulo $2^{S+1}$, which gives the displayed
endpoint progression with step $2\cdot3^d$.  Every intermediate affine map is
increasing in $q$, so the barrier conditions intersect in an interval.  Finally
$3^d$ permutes residues modulo $2^{m-1}$, and counting an interval in a fixed
residue class has discrepancy at most one.
\end{proof}

\begin{proposition}[Boundary-layer spectrum of the bad event]
\label{prop:bad_boundary_spectrum}
Let $K=2^M$, write an odd integer as $N=2z+1$, and put
\[
 \mathcal B_{d,M}=\{z\pmod K:S_d(2z+1)\le M\}.
\]
For a valuation word $w=(a_1,\ldots,a_d)$ write $S=|w|$ and define
$c_0=0$, $S_0=0$, and
$c_j=3c_{j-1}+2^{S_{j-1}}$, $S_j=S_{j-1}+a_j$.  Its exact odd-start
residue is
\[
 \rho_w\equiv(2^S-c_d)3^{-d}\pmod{2^{S+1}},
 \qquad r_w=(\rho_w-1)/2\pmod{2^S}.
\]
If
\[
 \beta_{d,M}(h)=\frac1K\sum_{z\pmod K}
 \mathbf1_{\mathcal B_{d,M}}(z)\exp(-2\pi i hz/K),
\]
then, for $h\ne0$ and $v=v_2(h)$,
\[
 \boxed{
 \beta_{d,M}(h)=
 \sum_{S=\max\{d,M-v\}}^M2^{-S}
 \sum_{\substack{w\in\mathbb N^d\\|w|=S}}
 \exp(-2\pi i h r_w/K).}
\]
Thus a frequency of valuation $v$ sees only the boundary layers
$M-v\le |w|\le M$; in particular, every odd frequency sees only $|w|=M$.
\end{proposition}
\begin{proof}
The exact valuation condition requires
$3^dN+c_d\equiv2^S\pmod{2^{S+1}}$, which gives the displayed residue.
In the $z$ coordinate the word cylinder is $z\equiv r_w\pmod{2^S}$.
Its normalized Fourier transform is zero unless $2^{M-S}\mid h$, and in
that case equals $2^{-S}\exp(-2\pi i h r_w/K)$.  Summing the disjoint exact
 word cylinders with $d\le S\le M$ proves the formula.
\end{proof}
\begin{lemma}[Top-layer offset injectivity]
\label{lem:top_layer_offset_injectivity}
Fix $d\le M$ and let $w,w'\in\mathbb N^d$ be distinct words of total
weight $M$.  Then their odd-coordinate residues satisfy
\[
 r_w\not\equiv r_{w'}\pmod{2^{M-1}}.
\]
Consequently the number of antipodal pairs on the top layer is zero:
$A_{\mathrm{ant}}=0$.  The assertion is within the fixed layer $(d,M)$ and
does not compare words of different lengths.
\end{lemma}
\begin{proof}
Write $S_j=a_1+\cdots+a_j$ and $S'_j=a'_1+\cdots+a'_j$, with
$S_0=S'_0=0$.  The affine constants have the expansion
\[
 c_w=\sum_{j=0}^{d-1}3^{d-1-j}2^{S_j},
 \qquad
 c_{w'}=\sum_{j=0}^{d-1}3^{d-1-j}2^{S'_j}.
\]
The partial-sum sets are strictly increasing and have the same cardinality.
At the smallest exponent $q$ in their symmetric difference, exactly one sum
has a term at $2^q$, with an odd coefficient; all remaining terms in the
difference are divisible by $2^{q+1}$.  Hence
$v_2(c_w-c_{w'})=q<M$, since every partial sum before the last is $<M$.
For a total weight-$M$ word,
\[
 \rho_w\equiv(2^M-c_w)3^{-d}\pmod{2^{M+1}},
 \qquad r_w=(\rho_w-1)/2.
\]
Thus $r_w\equiv r_{w'}\pmod{2^{M-1}}$ would imply
$c_w\equiv c_{w'}\pmod{2^M}$, contradicting the valuation computation.
\end{proof}

\begin{assumption}[Low-Frequency Cancellation Lemma; falsified]
\label{ass:low_freq_cancellation}
The unrestricted inter-word Fourier sum evaluated at any odd integer $h \neq 0$
satisfies
\[
|I(h)| = o(\sqrt{W_b}) \quad \text{as } W_b \to \infty.
\]
Consequently, the aggregate discrepancy satisfies
\[
\left|\sum_w \epsilon_w\right| \le \sum_{h=1}^{K} |\widehat{\mathbf{1}_{\mathcal{B}_b}}(h)| \cdot |I(h)| + O(\sqrt{W_b}),
\]
where the $O(\sqrt{W_b})$ term is the thermal noise from the remaining $W_b - K$ modes. If $|I(h)| = o(\sqrt{W_b})$ for all $h \le K$, then the mass-weighted testing error satisfies $\Delta_b \sim 2^{-M_b/2}$.
\end{assumption}

\begin{remark}[Falsification by unrestricted spectral diagnostics]
Exact DFS computation of the unrestricted Fourier spectrum at $M=22$ and
$M=24$ reveals persistent macroscopic peaks of magnitude $\sim 0.6$,
with no exponential decay. Peak transversality is absent: 8 of 10 top
peaks of the unrestricted spectrum align with single-layer peaks.
The noise tail (excluding structural peaks) retains peaks of magnitude
$\sim 0.43$, exceeding the $2^{-B/2}$ threshold by a factor of 885.
The Low-Frequency Cancellation Lemma is therefore falsified.
\end{remark}

\begin{observation}[Low-frequency structural cancellation; falsified]
\label{obs:low_freq_cancellation}
Fourier decomposition of the aggregate discrepancy $\sum_w \epsilon_w$ at
$M=22$ (65,535 modes) reveals a three-tier structure:
\[
\begin{array}{c|c|c}
\text{Modes} & \text{Fraction} & \text{Contribution to } |\sum\epsilon_w| \\ \hline
\text{Top 5} & 0.008\% & 58\% \text{ of real discrepancy} \\
\text{Top 20} & 0.03\% & 80\% \text{ of real discrepancy} \\
\text{Top 100} & 0.15\% & 95.5\% \text{ of real discrepancy} \\
\text{Top 1000} & 1.5\% & 99.6\% \text{ of real discrepancy} \\
\text{Remainder (64535)} & 98.5\% & 0.4\% \text{ (thermal noise } \approx 2.6\sqrt{W}\text{)}
\end{array}
\]
The dominant modes are the lowest global frequencies $h = \pm 1, \pm 2, \pm 3, \pm 4, \pm 5$. While this hierarchical mode distribution correctly describes the variance spread at small $M$, our exact unrestricted evaluations at $M=22, 24$ prove that the ratio $|\sum\epsilon_w|/\sqrt{\sum|\epsilon_w|}$ does not converge to zero. The structural component does not grow slower than the volume; the normalized Fourier peak retains $O(1)$ macroscopic magnitude.
\end{observation}

\begin{observation}[Hierarchical inter-bucket phase cancellation; falsified]
\label{obs:hierarchical_inter_bucket}
The hypothesis of coherent inter-bucket destructive interference is explicitly falsified. Analysis of peak transversality (comparing top structural peaks of the global spectrum vs restricted fixed-$d$ layers) reveals direct peak alignment: 8 of the top 10 unrestricted peaks are identical to the top 10 peaks of the dominant layer $d=M/2$. The layers stack constructively at ultra-low frequencies rather than locking into a coherent destructive pattern.
\end{observation}

\begin{lemma}[Spectral Transversality of Resonant Sets; falsified as a decay mechanism]
\label{lem:spectral_transversality}
Let $R_{\text{end}} = \{h : \|h 3^{d_{\text{end}}} / 2^M\| \le A/2^M\}$ be the set of resonant frequencies for the endpoint distribution, and $R_{\text{bad}} = \{h : \|h 3^{d_{\text{bad}}} / 2^M\| \le A/2^M\}$ be the resonant set for the bad-set boundary layer. 
While the arithmetic fact $3^{d_{\text{end}} - d_{\text{bad}}} \not\equiv 1 \pmod{2^M}$ holds, enforcing that the central peaks do not intersect, this separation does not imply smallness of the total error $\Delta_b$. The high-frequency noise cross-terms generate macroscopic peaks of magnitude $\sim 0.43$ independent of the resonance alignment, exceeding the required $2^{-B/2}$ threshold by a factor of 885.
\end{lemma}

\begin{remark}[Falsification of intra-layer anomalous cancellation]
Extensive exact numerical DP tracking over the bucket constraints (up to $M=10^5$) explicitly falsifies any hypothesis of anomalous intra-layer Fourier cancellation. Within any fixed, macroscopic slice of length $d \approx M/2$, the Fourier coefficients uniformly obey $\theta \approx 0.47 \approx 0.5$, mirroring an unstructured random walk. Anomalous suppression (such as $\theta \to 0.20$ observed in unrestricted phase sums) is exclusively a consequence of \emph{inter-layer} destructive interference, which is irrelevant for isolated buckets defined at fixed length $d$. Thus, the decay in $\Delta_b$ is not driven by the endpoint distribution becoming artificially smooth, but rather by its sharp peaks systematically avoiding the structural resonances of the bad set.
\end{remark}

\begin{theorem}[Conditional Forward Shadowing; hypothesis falsified]
\label{thm:conditional_shadowing}
Assume the Spectral Transversality of Resonant Sets (Lemma~\ref{lem:spectral_transversality}):
for each endpoint bucket $b$, the testing error against the bad-set indicator $\mathbf{1}_{\mathcal{B}_b}$ satisfies
\[
\Delta_b = \frac{|\sum_w \epsilon_w|}{Q_b} \sim C \cdot 2^{-M_b/2}.
\]
Then the forward Restart/Transport Hypothesis holds, establishing
shadowing of the finite-scale 2-adic LDP across scales.
\end{theorem}
\begin{remark}[Status]
The geometric decay condition $\Delta_b \sim C \cdot 2^{-M_b/2}$ is directly falsified by the macroscopic noise tail magnitude. The theorem lacks a valid hypothesis mechanism.
\end{remark}

\begin{observation}[Finite boundary-layer cancellation diagnostic]
At $\alpha=1.10$ exhaustive enumeration gives
\[
\begin{array}{c|ccc}
 B&16&18&20\\ \hline
 \max_b|\Delta_b|&0.004518&0.002422&0.001319\\
 \sum_bp_b|\Delta_b|&0.003506&0.001901&0.000887\\
 2^{B/2}\sum_bp_b|\Delta_b|&0.898&0.973&0.908
\end{array}
\]
The mass-weighted absolute error is strikingly consistent with
$2^{-B/2}$ over these three scales.  At $B=20$ the weighted absolute Fourier
majorant is $0.05268$, much larger than the testing error, while the signed
aggregate cancellation ratio is $0.00911$.  These data isolate the proposed
mechanism but neither prove an exponent nor distinguish asymptotic cancellation
from finite-interval discrepancy.
\end{observation}

\begin{observation}[Finite two-block restart diagnostic]
At $N_0=2^B$, $B=16,18,20$, and local exponents
$\alpha\in\{1.05,1.10,1.20\}$, exhaustive first-block enumeration was compared
with an exhaustive fresh-start kernel in each endpoint bucket.  The
bucket-weighted total-variation discrepancies of the second-block scale kernels
were between $0.005$ and $0.016$; the largest bucket discrepancy was below
$0.025$.  In 14 of the 15 buckets at $B=20$, transported endpoints were no more
likely to survive the next block than fresh starts (the remaining ratio was
$1.044$ in the smallest bucket).  These finite computations support the testing
form above but do not establish decay of the error as $B\to\infty$.
\end{observation}

\begin{observation}[Cocycle transient diagnostic; numerical]
For the non-autonomous cocycle $\mathcal{L}_{z,3^{d-1}\xi}\circ\cdots\circ\mathcal{L}_{z,\xi}$
with $\xi = 137/65536$ and $\xi \in \{1,3,1000\}/65536$, the per-step norm
$(\|\text{cocycle}_d\|)^{1/d}$ takes values $0.866, 0.927, 0.961, 0.979, 0.989$
at $d = 50, 100, 200, 400, 800$. The norm decreases monotonically toward $1.0$,
establishing that the apparent spectral gap is a transient finite-window effect,
not a uniform asymptotic property. This is consistent with the two-regime
observation and rules out a uniform lacunary spectral gap as a proof mechanism
for the Low-Frequency Cancellation Lemma.
\end{observation}

\begin{remark}[Reverse-chord encoding: superseded]
\label{rem:reverse_encoding_superseded}
The Reverse-chord Encoding Hypothesis was previously formulated as a
potential bridge to attach $\gamma_{\rm rev} \approx 0.993$ to the forward
exceptional set. Numerical experiments (E9) revealed a fundamental
information-theoretic barrier: the entropy of reverse chords with mean
shift $4/3$ is $h(4/3) \approx 1.077$ bits/step, while the target space
requires $\log_2 3 \approx 1.585$ bits/step. The resulting Hausdorff
dimension deficit ($D_{\rm rev} \approx 0.808 < D_{\rm fwd} \approx 0.950$)
makes 100\% coverage algebraically impossible. This route is closed.
\end{remark}

\begin{remark}[The path forward: no surviving analytic bridge]
With the Low-Frequency Cancellation Lemma (Assumption~\ref{ass:low_freq_cancellation})
and Spectral Transversality (Lemma~\ref{lem:spectral_transversality}) both
falsified as decay mechanisms by the unrestricted spectral diagnostics
(Observations~\ref{obs:low_freq_cancellation}, \ref{obs:hierarchical_inter_bucket}),
the forward descent route currently has \emph{no} candidate analytic bridge.
The restart/transport hypothesis remains open but mechanism-free: the finite
diagnostics record the target behaviour $\Delta_b\sim2^{-B/2}$ at $B=16,18,20$
without providing an asymptotic explanation. Any future route must bypass
Fourier $L^2/L^\infty$ norms entirely.
\end{remark}

\section{Caustic Scaling}

\begin{theorem}[C1: caustic scaling of confluence centers]
For the confirmed confluence centers spanning peaks $14 \le P \le 140$,
the affine law
\[
\mathrm{bits}(c) = a\,P + b,\qquad a = 0.4952,\quad b = 6.5567,\quad R^2 = 0.9723,
\]
holds, and the slope is statistically indistinguishable from $\tfrac12$
(the independent primary-sample fit gives $a=0.4983$, $R^2=0.965$).
Consequently
\[
\alpha_{\mathrm{eff}}(P)=\frac{\mathrm{bits}(c)}{P}=a+\frac{b}{P}\ \longrightarrow\ \frac12,
\]
with the finite-size intercept $b\approx6.56$ accounting for the inflated
small-peak values ($\alpha_{\mathrm{eff}}(14)\approx0.96$, Class~B mean
$\approx0.71$ matching $0.5+6.5/\langle P\rangle\approx0.70$, and
$\alpha_{\mathrm{eff}}(140)=0.536$).
\end{theorem}

\begin{remark}[Robustness and the 39-center refit]
The refit over all 39 centers including the five small Expedition~D centers
($0.6201\,P+2.2850$, $R^2=0.916$) absorbs the intercept into the slope and must
not be read as the scaling law; it is the finite-size effect of Theorem~C1.
\end{remark}

\subsection{The Caustic Balance: Detailed Balance at the Waist (Model C1)}

\begin{lemma}[Detailed balance of the conditioned reverse process (open)]
\label{lem:detail_balance}
Let $c$ be a confluence center of peak $P$, $b=\mathrm{bits}(c)$, and let
$\mathcal T_k(c)$ be the reverse tree of $c$ truncated to depth $k$, restricted to
preimages $y<2^P$ whose trajectory peaks at $P$. Let
$h_{\mathrm{rev}}:=\lim_{k\to\infty}\tfrac1k\log_2|\mathcal T_k(c)|$ be the asymptotic
entropy rate of the conditioned reverse tree, and
$r_{\mathrm{fwd}}:=(P-b)/d_{\mathrm{fwd}}$ the forward gain rate of the canonical path
$c\to P$. Then
\[
h_{\mathrm{rev}}=r_{\mathrm{fwd}},\qquad
\tfrac1k\log_2|\mathcal T_k|=h_{\mathrm{rev}}+o(1).
\]
\end{lemma}

\begin{theorem}[C1: the waist is the log-midpoint, conditional on
Lemma~\ref{lem:detail_balance}]
\label{thm:C1}
Under Lemma~\ref{lem:detail_balance}, the node of maximal coalescence satisfies
$\mathrm{bits}(c)=\tfrac12 P+O(1)$, i.e.\ $\alpha=\tfrac12$ asymptotically.
\end{theorem}

\begin{proof}[Proof sketch]
The forward funnel $c\to P$ spans height $P-b$ at log-volume rate $r_{\mathrm{fwd}}$;
the conditioned reverse tree spans height $b$ at entropy rate $h_{\mathrm{rev}}$. The
center is the unique node where the two cones' log-volume growth rates coincide
(Lemma). Since the two cones share the same rate and together span $P$, self-consistency
forces $b=P-b$ up to $O(1)$.
\end{proof}

\begin{remark}[Numerical confirmation and the Zone~2 artifact]
\texttt{balance\_check.py} confirms $|h_{\mathrm{rev}}-r_{\mathrm{fwd}}|\le0.04$ for the
primary centers (peaks 14, 19, 26, 32, 36). For Zone~2 (peak 140),
$r_{\mathrm{fwd}}=0.2510$ matches the vacuum rate $\log_2 3-4/3\approx0.251$ exactly; the
apparent $h_{\mathrm{rev}}\approx2.68$ is a transient of the first $\sim15$ layers (raw
branching $\approx2^{2.68}$). Over the full depth $d_{\mathrm{core}}=251$ the mean bit
growth is $65/251=0.259\approx r_{\mathrm{fwd}}$, restoring the balance. The lemma is
therefore a statement about the \emph{asymptotic} slope, not the initial layers.
\end{remark}
\subsection{Phases 4--5: Solitons as Tilted-Typical Paths, Centers as Dressed Vacua}

\begin{theorem}[S1: soliton dominance and tilted typicality]
\label{thm:S1}
\begin{enumerate}
\item[(i)] \textbf{Soliton dominance.} Let $f(P)$ be the fraction of the peak-$P$
flow that passes through the canonical confluence center. Then
$f(14)\approx0.387$, $f(16)\approx6\times10^{-3}$, and $f(P)\le10^{-3}$ for all
verified $P\ge22$. Hence asymptotically almost all peak-$P$ trajectories are
solitons (they avoid the center).
\item[(ii)] \textbf{Spectral concentration.} The shift density of solitons concentrates
in $S/d\in[1.21,1.23]\subset[1,\,1.33]$.
\item[(iii)] \textbf{Tilted typicality (conditional).} Under the valuation heuristic
(Tao, Heuristic~1.8), conditionally on attaining peak $P$ from a $b$-bit input, the
shift vector is a typical path of the Cram\'er-tilted $\mathrm{Geom}(2)$ measure with
tilt $s^*=s^*(\sigma)$, $\sigma=\log_2 3-(P-b)/d$; the concentration in (ii) is the
law of large numbers under this tilted measure, with tilted mean $\sigma^*\approx1.22$.
\end{enumerate}
\end{theorem}

\begin{remark}[Status of S1]
Items (i)--(ii) are computational theorems (\texttt{phase4\_solitons.py}). Item (iii)
is conditional on the valuation heuristic, rigorous in the regime $n\ll\log N$ by
Tao's Proposition~1.9. The window $[1,1.33]$ is the admissible chord window; the
observed $[1.21,1.23]$ is the tilted mean, not a new constant.
\end{remark}

\begin{theorem}[I1: core as dressed vacuum, sparse defect gas]
\label{thm:I1}
\begin{enumerate}
\item[(i)] \textbf{Gain/charge balance.} For a center $c$ of peak $P$ with core length
$d$, $\ \delta \;=\; d\big(\log_2 3-\tfrac43\big)-\big(P-\mathrm{bits}(c)\big)$.
\item[(ii)] \textbf{$O(1)$ charge.} Over the verified archipelago $P\in[14,50]$,
$|\delta|\le3.8$; for the giant caustic Zone~2 ($P=140$) $\delta=0.1589$.
In particular $\delta$ does not grow with $P$.
\item[(iii)] \textbf{Interpretation.} The core is the $(1,1,2)$ vacuum
$\xi=-29/11$ dressed by a sparse gas of topological charges (defects $a\ge3$);
(ii) states that no macroscopic charge accumulates, i.e.\ $\delta/d\to0$, so the
core is asymptotically the pure vacuum.
\end{enumerate}
\end{theorem}

\begin{remark}[Status of I1]
(i) is an identity (definition of $\delta$); (ii) is a computational theorem
(\texttt{phase5\_instantons.py}) over the finite verified range $P\le140$, so the
asymptotic claim $\delta/d\to0$ in (iii) is an extrapolation beyond the verified
range, stated as a conjecture. (iii) matches Theorem~T2 of the map paper.
\end{remark}

\begin{remark}[Generic vs.\ rigid: crystals cannot be grown]
Theorems~S1 and~I1 complete the picture. The \emph{generic} peak path is the
soliton, a typical path of the tilted measure (S1); the confluence center is a
\emph{rare rigid} object, a dressed vacuum with $O(1)$ charge (I1). Rigid objects
thus form a measure-zero subset of the tilted ensemble: they can be found but not
grown.
\end{remark}

\subsection{Macro-obstruction to reverse synthesis}

\begin{theorem}[Macro-obstruction to reverse synthesis]
\label{thm:macro_obstruction}
Let $c$ be a confluence center of peak $P$ with $\mathrm{bits}(c)=B$, and let $w$
be the shift vector of its canonical trajectory with total shift $S=|w|$.
Then $c$ is the unique $B$-bit integer in its $2$-adic cylinder modulo $2^{S}$.
By the caustic scaling (Theorem~C1), $B=\tfrac12 P+O(1)$, while the ascent chord
satisfies $S=\sigma(P-B)/(\log_2 3-\sigma)$ with $\sigma\approx1.33$, hence
$S\approx2.6P$ and the modular deficit $\Delta:=S-B$ obeys $\Delta\ge2P$ for all
sufficiently large $P$. Consequently the expected number of $B$-bit integers
realizing a prescribed macro-chord is $2^{-\Delta}=2^{-\Theta(P)}$, and reverse
synthesis by CRT has success probability $2^{-\Theta(P)}$.
\end{theorem}

\begin{corollary}[Existence is ontological, not algorithmic]
Macro-centers are not constructible; their existence is governed by the LDP
rates (Theorems~F1/R1) and is a measure-zero arithmetic accident. This is the
obstruction counterpart of ``crystals can be found, but cannot be grown.''
\end{corollary}

\begin{remark}[Validation]
All $15$ known centers (peaks $14$--$33$) satisfy $\Delta>0$
(e.g.\ peak $14$: $\Delta=25$; peak $24$: $\Delta=123$; peak $33$: $\Delta=92$),
confirming the obstruction concerns constructibility, not existence.
For peak $1400$: $B\approx699$, $S\approx3656$, $\Delta\approx2957$, expected count
$2^{-2957}\approx0$.
\end{remark}

\section{Front C and Fine-scale mixing}
\label{sec:front_c}

\begin{assumption}[Front C target: exponential white-point moment with intrinsic ceiling]
\label{lem:front_c2}
Fix $0<\varepsilon<1/100$ and put
$\eta(\varepsilon)=-\log\cos(\pi\varepsilon)$.  Let $N_W=N_W(n,\xi)$
be the number of white renewal positions, where a renewal at position $j$ means
$b_j=3$ and hence, conditionally,
\[
 L_j=L_{j-1}+3,\qquad L_{j-1}\sim\operatorname{NB}(2(j-1),1/2).
\]
The open Front--C target is that for every fixed $\varepsilon$ there are
$C(\varepsilon)<\infty$ and $c_{\rm res}(\varepsilon)>0$ such that
\[
 |S_\chi(n)|\ \le\ \mathbb E e^{-\eta N_W}
 \ \le\ C(\varepsilon)e^{-c_{\rm res}(\varepsilon)n}
\]
uniformly in $n$ and $3\nmid\xi$.  The constant must depend on $\varepsilon$;
in particular the data are consistent with
$c_{\rm res}(\varepsilon)=O(\varepsilon^2)$ for the cosine weight as
$\varepsilon\downarrow0$.  This estimate is not proved here.

Independently of the open upper bound, the $b_j=3$-only mechanism has the
intrinsic ceiling
\[
 \mathbb E e^{-\eta N_W}\ \ge\ (3/4)^{\lfloor n/2\rfloor}
 =e^{-\frac12\ln(4/3)n},
\]
so it cannot yield an exponent larger than
$\tfrac12\ln(4/3)\approx0.1438$ nats.  Earlier small-$n$ calibration
($n\le14$ and six frequencies) is evidence only and is superseded by the
large-$n$ two-regime diagnostics below.
\end{assumption}

\begin{remark}[Status of the ceiling and mechanism at worst frequencies]
The ceiling $(3/4)^{\lfloor n/2\rfloor}$ is a strict limitation of the white-point mechanism itself, not of the true Fourier decay magnitude $|S_\chi(n)|$. At the absolute worst-case resonant frequencies ($\xi = 2^{s^*}$), the purely white-point paths are almost completely dephased (e.g., $|A| \approx 2 \times 10^{-6}$ for paths with no renewal at $n=12$). The resonance is entirely carried by the fractal paths containing renewal states ($b_j = 3$), where black points perfectly cohere at $\xi = 2^{s^*}$ while the remainder of the renewal dynamics only partially dephases the signal. The limit mechanism is thus a complex variational compromise within the renewal paths (the free energy of a polymer), not the trivial ceiling.
\end{remark}

\begin{observation}[Two anti-concentration regimes; numerical]
\label{obs:two_ac_regimes}
For
\[
p_{n,j}(\xi;\varepsilon)=\mathbb P\!\left(\|\theta(j,L_{j-1}+3;\xi)\|
\le\varepsilon\mid b_j=3\right),\qquad
B_n(\xi;\varepsilon)=\frac1{\lfloor n/2\rfloor}\sum_{j\le n/2}p_{n,j}(\xi;\varepsilon),
\]
the exact negative-binomial computation and a $30{,}000$-path lower-tail
calibration reveal two finite-$n$ regimes at $\varepsilon=0.05$:
\[
\begin{array}{c|cc|cc}
 n & B_n(\xi_{\rm res}) & B_n(\xi_{\rm rnd})
   & -n^{-1}\log\mathbb E e^{-\eta N_W}(\xi_{\rm res})
   & -n^{-1}\log\mathbb E e^{-\eta N_W}(\xi_{\rm rnd})\\ \hline
 200&0.7964&0.1239&3.13\times10^{-4}&1.35\times10^{-3}\\
 400&0.8062&0.1119&2.98\times10^{-4}&1.37\times10^{-3}
\end{array}
\]
Here $\xi_{\rm res}$ maximizes $B_n$ among the tested low-discrete-log family
$\xi=2^s$, $0\le s<4n$, together with $500$ deterministic random units;
$\xi_{\rm rnd}$ is the worst member of that random sample.  Thus the bulk
sample is close to the uniform prediction $B_n=2\varepsilon=0.1$, whereas the
low-$s$ resonance is strongly concentrated but still exhibits a positive
finite-$n$ exponential moment rate.  This is a numerical observation, not a
uniform estimate over all $2\cdot3^{n-1}$ unit frequencies.
\end{observation}

\begin{observation}[Alignment of the two numerical resonances]
\label{obs:resonance_alignment}
Let $s_{\rm adv}$ maximize the Front--C black density over the tested low-$s$
window, and let $s_{\rm E6}$ minimize the polymer Fourier magnitude.  The two
distinct optimization problems satisfy numerically
\[
\begin{array}{c|rrrr}
 &\varepsilon=0.10&0.05&0.02&0.01\\ \hline
 n=200&s_{\rm E6}+2&s_{\rm E6}+2&s_{\rm E6}&s_{\rm E6}-1\\
 n=400&s_{\rm E6}+2&s_{\rm E6}+2&s_{\rm E6}&s_{\rm E6}-1
\end{array}
\]
with $s_{\rm E6}=310$ and $627$, respectively.  The nontrivial content is the
low discrete logarithm $s=O(n)$ and the $O(1)$ alignment, since $2$ generates
the entire unit group modulo $3^n$ and hence every unit is formally a power of
$2$.  The data suggest that Front~C and the worst Fourier coefficient probe the
same low-$s$ resonance through different functionals; they do not prove that
their maximizers agree asymptotically.
\end{observation}

\begin{assumption}[Two-regime Front--C target]
\label{ass:two_regime_front_c}
For every fixed sufficiently small $\varepsilon>0$ there exists
$c_{\rm res}(\varepsilon)>0$ such that the exponential-moment estimate in
Assumption~\ref{lem:front_c2} holds uniformly over all unit frequencies.
Moreover, outside an exceptional family of vanishing relative size, it holds
with a larger bulk rate $c_{\rm bulk}(\varepsilon)>c_{\rm res}(\varepsilon)$.
The computations suggest, only at $\varepsilon=0.05$,
$c_{\rm res}\approx3.0\times10^{-4}$ and
$c_{\rm bulk}\approx1.36\times10^{-3}$.  Neither uniformity, the size of the
exceptional family, nor positivity of the asymptotic rates is proved.
\end{assumption}

\section*{Lemma 0: exponential fine-scale mixing from exponential Fourier decay}

\begin{assumption}[T0$_{\mathrm{cert}}$]
There exist explicit $C_F,c>0$ such that for all $n\ge1$ and all
$\xi\in\mathbb Z/3^n\mathbb Z$ with $3\nmid\xi$,
$|\hat\mu_n(\xi)|=|\mathbb E\,e^{-2\pi i\xi\,\mathrm{Syrac}(\mathbb Z/3^n\mathbb Z)/3^n}|
\le C_Fe^{-cn}$.
Numerical status: certified pointwise only on finite ranges (the original exact
scan reached $n\le16$); the polymer recursion through $n=400$ gives a drifting
worst-case rate and suggests $c_\infty\approx0.053$--$0.058$ under the tested
fit families, but neither positivity nor the limiting value is certified (see
Observation~\ref{obs:two_regimes}).  The value
$\tfrac12\ln(4/3)\approx0.1438$ is the white-point mechanism ceiling, not the
actual worst-case rate.
\end{assumption}

\begin{lemma}[0a: linear-window typicality]
Let $(a_1,\dots,a_n)\equiv\mathrm{Geom}(2)^n$. For every $C>0$ there is
$b(C)>0$ such that the event
$|a[i,j]-2(j-i)|\le C\sqrt{n(j-i)}$ for all $1\le i\le j\le n$
has probability $\ge1-e^{-b(C)n}$.
\end{lemma}
\begin{proof}[Sketch] Bernstein for each interval: exponent
$\asymp C^2n(j-i)/(j-i)=C^2n$; union bound over $\le n^2$ intervals
preserves exponentiality. This replaces Tao's $\sqrt{(j-i)}\log n$ windows,
whose $\sim n^2\cdot n^{-cC_A}$ failure probability is only polynomial and
would cap the output at $(\log x)^{-K}$.
\end{proof}

\begin{lemma}[0: no additive gate]
Under T0$_{\mathrm{cert}}$, for all $10\le m\le n$,
$\mathrm{Osc}_{m,n}\le C'e^{-c_1m}$, where one may take explicitly
\[
c_1(c)=\frac{0.215\,c^2}{\ln2+c}\ >\ 0
\qquad\big(\text{feasible slack }\eta^*(c)=\tfrac{0.215c}{\ln2+c}\big).
\]
Numerically: $c_1(0.1438)\approx0.0053$, $c_1(0.25)\approx0.014$,
$c_1(0.55)\approx0.052$, $c_1(0.61)\approx0.061$. Constants are not optimized.
\end{lemma}
\begin{proof}[Proof skeleton]
Reconstruct Tao's \S6 with exponential bookkeeping. Regime $0.9n\le m\le n$
first (general $m$ by the lossless telescoping of Step~C). Lower the crossing
threshold to $L=n\log_2 3-\eta n$ (linear slack); by Lemma~0a the stopping time
satisfies $k\le0.7925n-\eta n/2+O(1)$ except $e^{-\Omega(n)}$. The tail factor
is $\le C_Fe^{-c\,\mathbf n'}$ with
$\mathbf n'=n-k-v_3(\xi)-1\ge m-k\ge(0.1075+\eta/2)n$. The Renyi factor with
the linear window costs $3^n\sum\mathbb P^2\le e^{\eta n}$ (Young's inequality
converts the $C\sqrt{nj}$ window into an $e^{O(n)}$ factor absorbed by the
slack, and Corollary~6.3-type injectivity survives since
$\sum_j3^{j-1}2^{l-a[1,j]}<3^n$ still closes). Assembling by Plancherel and
Cauchy--Schwarz:
$\mathrm{Osc}\le\exp\big(-c(0.1075+\tfrac\eta2)n+\tfrac\eta2n\ln2+O(\log n)\big)$.
The exponent $\varphi(\eta)=0.1075c-\tfrac\eta2(\ln2-c)$ is decreasing in
$\eta$ (since $c<\ln2$); the minimal feasible slack
$\eta^*=0.215c/(\ln2+c)$ (balancing the window constant against the slack)
gives the stated $c_1=\varphi(\eta^*)$.
\end{proof}

\begin{remark}[The gate of Path~B was an accounting artifact]
There is no additive threshold on $c$: every positive Fourier exponent yields
a positive mixing exponent. The price is multiplicative,
$\kappa(c)=c_1/c=0.215c/(\ln2+c)\in(0,0.215)$, not a gate at $0.494$.
\end{remark}

\begin{corollary}[Conditional stabilization calculation]
Assume T0$_{\mathrm{cert}}$ with a genuine asymptotic exponent $c>0$ and assume
the remaining hypotheses in the fine-scale gluing argument.  Then
$d_{TV}(\mathrm{Pass}_x(\mathbf N_{x\alpha}),\mathrm{Pass}_x(\mathbf N_{x\alpha^2}))
\le Cx^{-\gamma_{\mathrm{Tao}}}$ with
$\gamma_{\mathrm{Tao}}=c_1(c)(\alpha-1)/100$.  This is a conditional conversion
formula only.  The number $0.1438$ is the ceiling of the white-point mechanism,
not a certified Fourier exponent, and the finite-$n$ value near $0.25$ cannot
be substituted as an asymptotic rate.
\end{corollary}

\subsection{Two Fourier regimes and the worst-case barrier}

\begin{observation}[Two Fourier decay regimes; numerical]
\label{obs:two_regimes}
The Syracuse measure exhibits two numerically distinct Fourier behaviours:
\begin{enumerate}
\item[(i)] \textbf{Bulk ($L^2$).} The collision dimension satisfies $D_2\to1$ and
the root-mean-square decay rate is $c_{\mathrm{avg}}\approx0.61$ (nats),
reflecting chaotic mixing of the typical mass.  The independent Front--C
sample in Observation~\ref{obs:two_ac_regimes} likewise has black density close
to $2\varepsilon$ and a substantially faster exponential-moment rate than the
resonant sample.
\item[(ii)] \textbf{Worst case ($L^\infty$).} The supremum
$\sup_{3\nmid\xi}|\hat\mu_n(\xi)|$ is attained in the computations at a
low-discrete-log frequency $\xi=2^{s^*}$ with $s^*/n\to\log_2 3$.  The
worst-case rate $c_\infty$ satisfies $c_\infty\le0.054$ in the available fits;
the data do not exclude $c_\infty=0$ (subexponential decay).  The relevant
exceptional set is the truncated low-exponent family
$\{2^s:0\le s\le Cn\}$, whose cardinality is $O(n)$ compared with
$\varphi(3^n)=2\cdot3^{n-1}$; the full orbit $\{2^s\bmod3^n\}$ is the entire
unit group and must not be called $L^2$-negligible.
\end{enumerate}
\end{observation}

\begin{remark}[Status of Observation~\ref{obs:two_regimes}]
Both items are numerical.  The low-$s$ identification lacks a theoretical
explanation, the search at large $n$ does not exhaust all unit frequencies,
and $c_\infty$ is not determined.  Front--C alignment
(Observation~\ref{obs:resonance_alignment}) is independent supporting evidence,
but it does not turn either numerical maximizer into a theorem.
\end{remark}

\begin{theorem}[Worst-case barrier; cf.\ Tao, Remark 1.18]
\label{thm:barrier}
For $\xi$ not divisible by $3$,
$\ |\mathbb E e^{-2\pi i\xi\,\mathrm{Syrac}(\mathbb Z/3^n\mathbb Z)/3^n}|
\le\mathrm{Osc}_{n-1,n}.$
Hence any exponential fine-scale bound $\mathrm{Osc}_{m,n}\le Ce^{-cm}$ at
$m=n-1$ is limited by the worst-case Fourier rate $c_\infty$: if $c_\infty=0$,
Tao's superpolynomial bound cannot be upgraded to exponential.
\end{theorem}

\begin{remark}[$L^2$ optics cannot bypass the barrier]
Although the bulk $L^2$ decay is fast, the finest-scale oscillation
$\mathrm{Osc}_{n-1,n}$ is bounded \emph{below} by the worst-case Fourier
coefficient (Theorem~\ref{thm:barrier}). Averaging over frequencies therefore
cannot yield an exponential bound when the worst-case decay is subexponential;
the barrier is intrinsic to the equivalence of fine-scale mixing and pointwise
Fourier decay.
\end{remark}

\subsection{Honest status of the programme}

\begin{theorem}[Conditional reverse-geometry consequence]
\label{thm:final}
If the reverse-chord encoding were valid, the reverse-growth calculation
would assign the candidate exponent
\[
 \gamma_{\rm rev}
 =\frac{I_{\rm rev}(1.33)}{\log_2 3-1.33}\approx0.993
\]
to the forward exceptional set. However, as established in
Remark~\ref{rem:reverse_encoding_superseded}, the encoding suffers an intrinsic
entropy deficit. Therefore, this number is strictly a characteristic exponent of
reverse growth geometry and cannot bound $E_{N_0}$.
\end{theorem}

\begin{remark}[What remains open]
\begin{enumerate}
\item[(i)] The restart/transport hypothesis across scales is open and currently
mechanism-free. Every candidate mechanism --- pointwise Fourier decay,
low-frequency/aggregate cancellation, spectral transversality --- has been
falsified by the unrestricted spectral diagnostics
(Assumption~\ref{ass:low_freq_cancellation}, Lemma~\ref{lem:spectral_transversality}).
\item[(ii)] The worst-case Fourier rate $c_\infty$ is not determined. If
$c_\infty=0$, Tao's architecture yields a superpolynomial but not an
exponential exceptional-set bound.
\item[(iii)] The full Collatz conjecture $\mathrm{Col}_{\min}(N)=1\ \forall N$
is a pointwise statement. It is not implied by any density bound and remains
beyond the statistical methods developed here.
\end{enumerate}
\end{remark}

\section{Conclusion: Local Descent, Reverse Geometry, and the Open Bridges}
\label{sec:conclusion}

The reverse thermodynamic calculation gives the dimensionless geometric ratio
\[
 \gamma_{\rm rev}
 =\frac{I_{\rm rev}(1.33)}{\log_2 3-1.33}
 \approx\frac{0.2532}{0.25496}\approx0.993.
\]
\paragraph{The closure of the Fourier route.}
The exact boundary-layer spectrum formula and the affine word structure are mathematically sound and robustly validated. However, the analytic path to Forward Shadowing via Fourier-based max bounds or L2 cancellation is conclusively closed. The unrestricted spectrum preserves macroscopic structural peaks, and the high-frequency ``noise'' retains $O(1)$ magnitude, explicitly failing the required $2^{-B/2}$ decay threshold. Consequently, neither pointwise Fourier suppression nor Spectral Transversality provides the global error cancellation necessary for an asymptotic proof. The numerical $2^{-B/2}$ decay observed in specific finite-scale bad sets is a local phenomenon reflecting specific topological properties of those sets, rather than a universal property of the endpoint measure. Any path forward must circumvent these Fourier obstructions entirely.

\paragraph{Computational frontier.}
The numerical substitutions at $N_0=2^{68}$ previously reported from
$K N_0^{-0.993}$ are therefore conditional on the unproved reverse-chord
encoding, not consequences of the reverse LDP alone.  They should be read only
as a scale illustration under that stronger hypothesis, not as a residual
density established by the present work.

\paragraph{Honest status and limits of the proof.}
It is crucial to clarify the rigorous boundaries of our results:
\begin{enumerate}
\item \textbf{Rigorous 3-adic Thermodynamics:} Steps A and B prove the properties of the tilted 3-adic model strictly. The spectral gap, the Doob $h$-transformed LLN, and the variance of the conditioned process are rigorous theorems about the 3-adic topology, confirmed by numerical simulation up to $n=7$.
\item \textbf{Two Open Bridges:} The forward route requires restart/transport of the endpoint distribution between local $2$-adic LDP blocks. The reverse value $\gamma_{\rm rev}\approx0.993$ additionally requires a combinatorial encoding of every forward exceptional integer by a reverse $4/3$-shift chord. Neither bridge is proved.
\item \textbf{The Worst-Case Barrier:} The Fourier route (Tao's architecture) unconditionally relies on the fine-scale mixing equivalence. If the worst-case Fourier decay is subexponential ($c_\infty = 0$), then the architecture yields only a superpolynomial exceptional-set bound, not an exponential one. The $L^2$ optics cannot bypass this barrier at the finest scale $m=n-1$.
\end{enumerate}

In conclusion, the project currently contains a rigorous finite-scale $2$-adic counting bound, conditional thermodynamic models, and numerical Fourier and reverse-growth observations.  The value $\gamma_{\rm rev}\approx0.993$ is a reverse-geometric characteristic, not an established exceptional-set exponent.  Establishing a forward bound requires restart/transport; attaching the reverse exponent requires reverse-chord encoding.  The full Collatz conjecture remains beyond these density methods.

\end{document}
